\documentclass[11pt,a4paper]{article}

% --- Encodage / langue ---
\usepackage[utf8]{inputenc}
\usepackage[T1]{fontenc}
\usepackage[french]{babel}
\usepackage{lmodern}

% --- Mise en page compacte (tient en 1 page hors illustrations) ---
\usepackage[a4paper,margin=2cm]{geometry}
\usepackage{microtype}
\usepackage{enumitem}
\setlist{nosep,leftmargin=1.2em}
\usepackage{titlesec}
\titlespacing*{\section}{0pt}{6pt}{2pt}
\titlespacing*{\subsection}{0pt}{4pt}{1pt}
\titleformat{\section}{\normalfont\large\bfseries}{\thesection.}{0.5em}{}
\titleformat{\subsection}{\normalfont\bfseries}{}{0em}{}
\setlength{\parindent}{0pt}
\setlength{\parskip}{2pt}

\usepackage[backend=biber,style=numeric-comp,sorting=none,maxbibnames=3]{biblatex}
\addbibresource{references.bib}
\renewcommand*{\bibfont}{\footnotesize}

\usepackage{hyperref}
\hypersetup{hidelinks}

\title{\vspace{-1.2cm}\textbf{État de l'art --- Analyse automatique de partitions musicales}}

\author{\normalsize Thibault Bourdoiseau \& Léo Li --- SAE TIFO 1 / TIFO 2 / MLRF}
\date{\vspace{-8pt}\small \today}

\begin{document}
\maketitle
\vspace{-12pt}

\section{Problème et objectif}
La reconnaissance optique de partitions (OMR) consiste à transformer la photo ou le scan d'une partition
imprimée en une représentation que l'ordinateur peut exploiter, c'est-à-dire la liste des notes avec leur
hauteur et leur durée, les clefs et les silences, avec éventuellement un export au format MIDI. La
littérature de référence \cite{rebelo2012,bainbridge2001,calvo2020} décrit pour cela une chaîne de
traitement bien établie. Cette chaîne peut être réalisée entièrement avec les notions vues en TIFO 1,
TIFO 2 et MLRF, et donc sans recourir aux réseaux de neurones.

\section{Chaîne de traitement classique}

\subsection{1. Pré-traitement et binarisation}
On convertit d'abord l'image en niveaux de gris, puis on atténue le bruit à l'aide d'un filtre médian.
On sépare ensuite l'encre du papier en binarisant l'image : le seuil est calculé automatiquement par la
méthode d'Otsu \cite{otsu1979}, que l'on adapte localement lorsque l'éclairage n'est pas uniforme, comme
sur une photo prise au téléphone. Une ouverture morphologique élimine enfin les petits points de bruit
qui subsistent.

\subsection{2. Détection et suppression des portées}
C'est l'étape la plus délicate et la plus étudiée de l'OMR \cite{dalitz2008}, car les lignes de la portée
traversent les notes et gênent leur reconnaissance. Plusieurs approches sont possibles. On peut compter
les pixels noirs présents sur chaque ligne de l'image : les cinq lignes d'une portée ressortent alors sous
la forme de cinq pics réguliers qui les localisent, ce qui revient à utiliser un histogramme. On peut aussi
mesurer la longueur des suites de pixels noirs et blancs pour estimer l'épaisseur des lignes et l'écart
entre elles, deux valeurs qui serviront à régler toutes les étapes suivantes. Pour effacer les portées, on
applique une ouverture morphologique avec un élément structurant horizontal, puis on la soustrait à l'image
afin de retirer les lignes tout en gardant les notes. Enfin, lorsque les portées sont légèrement courbées
ou inclinées, la transformée de Hough ou la méthode des chemins stables \cite{cardoso2009} donnent de
meilleurs résultats.

\subsection{3. Segmentation des symboles}
Une fois les portées retirées, on regroupe les pixels connectés entre eux pour isoler chaque symbole.
On distingue ensuite les têtes de notes, les hampes et les barres de ligature grâce à des opérations de
morphologie mathématique et à l'analyse de la répartition des pixels.

\subsection{4. Détection et classification des symboles}
Pour repérer les têtes de notes, on peut comparer l'image à un modèle de référence, ou utiliser une
ouverture morphologique avec un élément structurant de forme ovale. Le fait que la tête soit pleine ou
vide donne déjà une indication sur la durée de la note. Pour les clefs, les silences et les altérations,
on décrit chaque forme par quelques mesures caractéristiques comme sa surface, son allongement ou ses
moments de Hu, puis on confie le classement à un algorithme simple tel que les $k$ plus proches voisins
ou une machine à vecteurs de support (SVM), vus en MLRF.

\subsection{5. Décodage musical et export}
La hauteur d'une note se déduit de la position verticale de sa tête par rapport aux lignes de la portée.
Sa durée se lit à partir du type de tête, pleine ou vide, et de la présence d'une hampe, de ligatures ou
de points. On reconstitue alors la suite des notes, que l'on peut exporter au format MIDI à l'aide d'une
bibliothèque comme \texttt{music21} \cite{cuthbert2010}.
\newpage
\printbibliography[heading=bibnumbered,title={Références}]

\end{document}
